\chapter{Endocrine Disorders and Diseases}
\index{Endocrine Disorders and Diseases}

\section*{Definition and Overview}
Endocrine disorders are conditions that affect the endocrine system, the network of glands that produce and release hormones. Hormones are chemical messengers that travel through the blood to regulate metabolism, growth, reproduction, mood, and many other body functions. The main glands are the pituitary, thyroid, adrenal glands, pancreas, ovaries, and testes. Endocrine disorders occur when a gland produces too much or too little of a hormone, or when the body does not respond to hormones normally. Common examples include diabetes, thyroid disease, and adrenal disorders.

\section*{Epidemiology}
Endocrine disorders are among the most common chronic conditions. Diabetes affects more than one in twenty people, thyroid disease affects a significant proportion of the population, and conditions such as polycystic ovary syndrome and osteoporosis are widespread. Hormonal disorders can affect people of any age, from congenital conditions in childhood to age-related changes in later life, and they account for a large share of GP and specialist care.

\section*{Aetiology and Pathophysiology}
Endocrine disorders arise from disruption at any level of the hormonal control system.
\begin{itemize}
    \item \textbf{Hypofunction:} a gland produces too little hormone, as in hypothyroidism, adrenal insufficiency, and type 1 diabetes
    \item \textbf{Hyperfunction:} a gland produces too much hormone, as in hyperthyroidism, Cushing's syndrome, and acromegaly
    \item \textbf{Hormone resistance:} the body does not respond normally to a hormone, as in type 2 diabetes
    \item \textbf{Feedback control:} the hypothalamus and pituitary regulate the glands, and problems anywhere in this system disrupt hormone levels
    \item \textbf{Causes:} autoimmune attack, tumours, genetic conditions, infection, injury, and medicines
    \item \textbf{Effects:} because hormones act throughout the body, the effects of endocrine disease are widespread, affecting metabolism, heart, bones, mood, and reproduction
\end{itemize}

\section*{Clinical Features}
The features depend on which hormone is affected and whether levels are high or low.
\begin{itemize}
    \item Weight change, fatigue, and temperature intolerance
    \item Changes in heart rate and blood pressure
    \item Changes in skin, hair, and sweating
    \item Mood changes, anxiety, and depression
    \item Menstrual irregularity and fertility problems
    \item Muscle weakness and bone pain
    \item Increased thirst and urination with diabetes
    \item Growth problems in children
    \item Features specific to each gland and hormone
\end{itemize}

\section*{Differential Diagnosis}
The symptoms of endocrine disease overlap with many other conditions.
\begin{itemize}
    \item Chronic fatigue, depression, and anxiety
    \item Malnutrition and eating disorders
    \item Chronic infections and inflammatory disease
    \item Cardiovascular and kidney disease
    \item Normal ageing, which can be mistaken for hormone deficiency
\end{itemize}

\section*{When to Seek Medical Care}
People with persistent fatigue, weight change, palpitations, or other hormonal symptoms should see a GP, who can arrange testing and refer to an endocrinologist. Some endocrine emergencies, such as diabetic ketoacidosis and adrenal crisis, require urgent care.

\textbf{Contact your local emergency services for:}
\begin{itemize}
    \item Very high blood sugar with vomiting, confusion, or rapid breathing
    \item Severe weakness, collapse, or vomiting in a person with adrenal disease
    \item Palpitations with chest pain or fainting
    \item Severe headache with visual changes, which may indicate a pituitary problem
\end{itemize}

\section*{Diagnosis and Investigations}
Diagnosis is based on hormone tests, often with stimulation or suppression testing.
\begin{itemize}
    \item \textbf{Blood tests:} hormone levels, including thyroid, cortisol, and glucose
    \item \textbf{Glucose tests:} fasting glucose and oral glucose tolerance tests for diabetes
    \item \textbf{Urine tests:} for hormone breakdown products
    \item \textbf{Imaging:} ultrasound, CT, or MRI of the glands
    \item \textbf{Stimulation tests:} assessing how glands respond to stimulation
    \item \textbf{Referral:} to an endocrinologist for complex cases
\end{itemize}

\section*{Management}
Management depends on the specific disorder.
\begin{itemize}
    \item \textbf{Hormone replacement:} for underactive glands, such as thyroxine for hypothyroidism and insulin for type 1 diabetes
    \item \textbf{Reducing hormone levels:} medicines, surgery, or radiotherapy for overactive glands and tumours
    \item \textbf{Medicines to block hormone effects:} for hormone resistance and overactivity
    \item \textbf{Lifestyle measures:} diet, exercise, and weight management, central to diabetes and metabolic disease
    \item \textbf{Monitoring:} regular hormone and complication checks
    \item \textbf{Surgery:} removal of overactive or cancerous glands
\end{itemize}

\section*{Complications}
\begin{itemize}
    \item Complications of diabetes, including heart, eye, kidney, and nerve disease
    \item Bone disease and fractures from hormonal imbalance
    \item Cardiovascular disease from thyroid and adrenal disorders
    \item Infertility and reproductive problems
    \item Growth and development problems in children
    \item Endocrine emergencies, including adrenal crisis and thyroid storm
\end{itemize}

\section*{Prognosis and Outcomes}
Most endocrine disorders are highly treatable, and with correct diagnosis and management, people live full and healthy lives. Diabetes, thyroid disease, and adrenal conditions are all managed effectively with modern treatment. The key to good outcomes is early diagnosis, appropriate specialist care, and consistent follow-up.

\section*{Prevention and Self-Care}
\begin{itemize}
    \item Attend regular health checks, including blood glucose and thyroid testing where indicated
    \item Take hormone medicines exactly as prescribed and never stop them abruptly
    \item Maintain a healthy diet and regular activity
    \item Report new or worsening symptoms promptly
    \item Carry medical identification if you have a condition such as adrenal insufficiency or diabetes
    \item Attend scheduled specialist reviews
\end{itemize}

\section*{Sources and Further Reading}
The following reputable sources provide further information for healthcare students and patients seeking reliable, current advice about endocrine disorders.
\begin{itemize}
    \item Endocrine Society of Australia. \textit{Endocrine disorders information.} https://www.endocrinesociety.org.au/
    \item Healthdirect Australia. \textit{Hormone and gland disorders.} https://www.healthdirect.gov.au/
    \item NHS. \textit{Hormonal and endocrine conditions.} https://www.nhs.uk/conditions/
    \item Mayo Clinic. \textit{Endocrine system disorders.} https://www.mayoclinic.org/diseases-conditions/
    \item MedlinePlus. \textit{Endocrine diseases.} https://medlineplus.gov/endocrinediseases.html
    \item MSD Manual. \textit{Endocrine and metabolic disorders.} https://www.msdmanuals.com/
\end{itemize}
